% Figure: isothermal pressure-composition (P-x-y) diagram for an ideal
% binary mixture obeying Raoult's law at fixed temperature. The straight
% upper line is the bubble curve, P = x1 P1sat + (1-x1) P2sat, linear in
% the liquid mole fraction x1. The curved lower line is the dew curve,
% P = 1/(y1/P1sat + (1-y1)/P2sat), plotted against the vapour mole
% fraction y1. Between the two lines the mixture is two-phase; a horizontal
% tie line at a given P connects the liquid composition x1 (on the bubble
% line) to the vapour composition y1 (on the dew line). Component 1 is the
% more volatile (P1sat = 1.0), component 2 the less volatile (P2sat = 0.4),
% in arbitrary pressure units scaled so P1sat = 1.
\begin{figure}[htbp]
\centering
\begin{tikzpicture}
\begin{axis}[
  width=12cm, height=8.5cm,
  xlabel={mole fraction of the volatile component $x_1$, $y_1$},
  ylabel={pressure $p/p_1^{\text{sat}}$},
  xmin=0, xmax=1, ymin=0.35, ymax=1.05,
  axis lines=left,
  tick label style={font=\scriptsize},
  label style={font=\small},
  legend pos=north west,
  legend style={font=\scriptsize},
  clip=false,
]
% Bubble line: liquid composition x1 -> P
\addplot[engblue, very thick] coordinates {
  (0.00,0.4000) (0.10,0.4600) (0.20,0.5200) (0.30,0.5800) (0.40,0.6400)
  (0.50,0.7000) (0.60,0.7600) (0.70,0.8200) (0.80,0.8800) (0.90,0.9400)
  (1.00,1.0000)
};
\addlegendentry{bubble line (liquid $x_1$)}
% Dew line: vapour composition y1 -> P
\addplot[engred, very thick] coordinates {
  (0.00,0.4000) (0.10,0.4255) (0.20,0.4545) (0.30,0.4878) (0.40,0.5263)
  (0.50,0.5714) (0.60,0.6250) (0.70,0.6897) (0.80,0.7692) (0.90,0.8696)
  (1.00,1.0000)
};
\addlegendentry{dew line (vapour $y_1$)}
% A tie line at P = 0.70: bubble x1 = 0.50, dew y1 solves 0.70 = 1/(y/1+(1-y)/0.4)
% -> 1/0.70 = 1.4286 = y + 2.5(1-y) = 2.5 - 1.5y -> y = (2.5-1.4286)/1.5 = 0.714
\addplot[enggray, thick, dashed] coordinates {(0.50,0.70) (0.714,0.70)};
\addplot[only marks, mark=*, engblack, mark size=1.8pt] coordinates {(0.50,0.70) (0.714,0.70)};
\node[anchor=north, font=\scriptsize] at (axis cs:0.50,0.685) {$x_1$};
\node[anchor=north, font=\scriptsize] at (axis cs:0.714,0.685) {$y_1$};
% Region labels
\node[font=\small\itshape] at (axis cs:0.28,0.95) {liquid};
\node[font=\small\itshape] at (axis cs:0.72,0.44) {vapour};
\node[font=\scriptsize] at (axis cs:0.52,0.80) {two-phase};
\end{axis}
\end{tikzpicture}
\caption[Isothermal $P$-$x$-$y$ diagram for an ideal binary]{Pressure
against composition for an ideal binary mixture at fixed temperature. The
straight bubble line gives the total pressure of a liquid of composition
$x_1$ (Raoult's law makes it linear); the curved dew line gives the pressure
at which a vapour of composition $y_1$ begins to condense. A horizontal tie
line at a chosen pressure joins the liquid composition on the bubble line to
the vapour composition on the dew line: at $p/p_1^{\text{sat}} = 0.70$ a
liquid at $x_1 = 0.50$ coexists with a vapour enriched to $y_1 = 0.71$ in
the volatile component. The vapour is always richer than the liquid in the
more volatile species, which is the separation that distillation exploits.
A flash calculation reads a feed of overall composition between the two
lines and splits it into these two coexisting phases.}
\label{fig:vol04ch02:binary-pxy}
\end{figure}
